% !TEX ROOT = ../Thesis.tex %-
%-> 中文摘要
%-
\chapter{摘\quad 要}\chaptermark{摘\quad 要}
\linespread{1.5}
\zihao{-4} 多元时间序列预测是刻画复杂动态系统演化规律的重要技术手段，在能源调度、交通管理、金融风控和气象预报等领域具有广泛应用价值。近年来，深度学习方法在非线性关系建模和长程依赖刻画方面取得了显著进展，但现有方法在多元长程预测中仍面临两个突出挑战：一是仅依赖时域建模时，难以同时兼顾局部突变、多尺度波动与长期周期结构；二是大语言模型所蕴含的语义先验与数值时间序列之间存在模态差异，直接融合容易带来对齐困难与额外计算开销。围绕上述问题，本文面向多元长程时间序列预测任务，研究如何在冻结大语言模型主体的前提下，将语义先验与时频域增强机制有效结合，以提升模型对长期趋势、多尺度波动和跨变量关联的刻画能力。

本文首先构建了一个统一预测骨架，将时域建模、基础频域表示和外部冻结语言模型生成的变量级语义嵌入纳入同一框架，并通过轻量适配、步长感知融合和跨模态注意力对齐实现数值特征与语义先验的协同建模。在此基础上，本文进一步提出两种频域增强方法。第一，提出融合小波包分解的局部多尺度时频增强方法，通过可学习小波包分解、子带级编码、频带交互与注意力池化，更细致地刻画局部波动、突变片段和多尺度周期结构。第二，提出融合频域前馈网络的全局频域增强方法，通过复数频谱建模、频域稀疏化和步长感知门控融合，强化模型对全局周期模式、主导频率结构和长期趋势信息的建模能力。

本文在多个公开多元时间序列基准数据集（ETT、Weather、ILI、Exchange 等）上进行了系统实验。结果表明，所提出的两种方法均能在统一骨架基础上提升长程预测性能，并在复杂波动场景和长预测步长设置下表现出较为一致的改进趋势。其中，融合小波包分解的方法能够更细致地刻画局部异常、高频扰动和多尺度波动结构；融合频域前馈网络的方法则在全局周期模式建模、长期趋势保持和整体预测稳定性方面表现出较好的能力。总体来看，两种频域增强方法在所选公开基准与代表性数据集上均取得了较好的综合效果，并体现出较强的互补性。

研究结果表明，在统一骨架下融合冻结大语言模型语义先验与频域增强策略，能够为多元长程时间序列预测提供一种具有较好扩展性和结构可解释性的技术路径。

\keywords{多元时间序列预测; 大语言模型; 小波包分解; 频域前馈网络; 时频增强}
%-
%-> 英文摘要
%-
\chapter{Abstract}\chaptermark{Abstract} Multivariate time series forecasting is a key technique for modeling the evolution of complex dynamic systems and is widely used in energy scheduling, traffic management, financial risk control, and weather forecasting. Although recent deep learning methods have significantly improved nonlinear modeling and long-range dependency learning, long-horizon multivariate forecasting still faces two major challenges: time-domain models often struggle to jointly capture local abrupt changes, multi-scale fluctuations, and long-term periodic structures, while the semantic priors contained in large language models are difficult to integrate directly with numerical time series because of the modality gap. To address these issues, this thesis investigates how frozen large-language-model semantic priors can be effectively combined with time-frequency enhancement mechanisms for long-horizon multivariate forecasting.

This thesis first builds a unified forecasting backbone. It integrates time-domain modeling, basic frequency-domain representations, and variable-level semantic embeddings generated by an external frozen language model. Lightweight adaptation, horizon-aware fusion, and cross-modal attention alignment are then introduced to promote collaborative modeling of numerical features and semantic priors. On this basis, two frequency-enhanced methods are further proposed. The first is a local multi-scale time-frequency enhancement method based on wavelet packet decomposition, which uses learnable wavelet packet decomposition, subband-level encoding, band interaction, and attention pooling to better characterize local fluctuations, abrupt segments, and multi-scale periodic patterns. The second is a global frequency enhancement method based on a frequency-domain feed-forward network, which employs complex-spectrum modeling, frequency sparsification, and horizon-aware gated fusion to strengthen the modeling of global periodic patterns, dominant spectral structures, and long-term trend information.

Extensive experiments on multiple public datasets, including ETT, Weather, ILI, and Exchange, show that both methods improve long-horizon forecasting performance over the unified backbone and exhibit relatively consistent gains under volatile scenarios and long prediction horizons. The wavelet-packet-based method provides finer modeling of local anomalies, high-frequency disturbances, and multi-scale fluctuations, while the frequency-domain-feed-forward-network-based method shows strong ability in modeling global periodic patterns, preserving long-term trends, and maintaining stable overall forecasting quality. Overall, both frequency-enhanced methods achieve solid comprehensive performance on the selected public benchmarks and representative datasets, and demonstrate clear complementarity.

These results demonstrate that combining frozen large-language-model semantic priors with local multi-scale and global frequency enhancement provides a scalable and structurally interpretable technical path for multivariate long-horizon time series forecasting.

\englishkeywords{multivariate time series forecasting; large language model; wavelet packet decomposition; frequency-domain feed-forward network; time-frequency enhancement}
